\section{Ingeniería de Detalles: viabilidad mecánica del prototipo}
Se utiliza exclusivamente el escenario acoplado M8 de 4 mm de biomasa:
presión de Holm 2,25867 MPa, coeficiente efectivo de torque 0,1,
fuerza máxima 918,00 N y torque máximo 2,7804 N\,m. El accionamiento es
manual, con manivela de 165 mm; el ABS es material de diseño de carcasa.
El modelo M8 de doble empotramiento y sus conclusiones se conservan, sin
repetir la comparación de diámetros ni el estudio de concentradores.
EP2 se emplea como respaldo de equilibrio, resistencia nominal, vida de
rodamientos y elasticidad tridimensional\cite{equipo_entrega2_2026}.

\paragraph{Árbol principal.}
Para la sección nominal circular de 12,7 mm se obtienen
$\sigma_{\rm axial}=7,25$ MPa, $\tau_T=6,91$ MPa y
$\sigma_{\rm VM}=14,00$ MPa en el mismo estado de carga.
Incluyendo la flexión manual y una envolvente correlacionada del pequeño
desequilibrio de compactación, la demanda nominal es 20,31 MPa.
Con $E=210$ GPa, la compresión sobre 200 mm es 0,00690 mm y el giro
torsional 0,1545 grados. Con dos guías radiales ideales en las placas CAD,
la flecha transversal estimada es 0,0183 mm. Estos valores corresponden
al diámetro exterior liso de referencia: el producto descrito en EP2 es
hilo corrido de media pulgada. Para cerrar su resistencia faltan el paso,
diámetro de raíz y grado del acero; no se asigna un factor de seguridad
al producto desconocido. El apoyo axial debe transmitir aproximadamente
0,924 kN al bastidor, incluyendo pesos de referencia de matriz y árbol;
no aparece un rodamiento axial instalado en el STEP.

\paragraph{Rodamientos de los rodillos.}
Con las reacciones correlacionadas y los pesos de los rodillos, la carga
radial máxima es 272,47 N. Para la referencia EZO 608ZZ,
$C=3300$ N y $C_0=1370$ N\cite{ezo_608zz_2026}, resulta
$S_{0,\min}=5,03$: capacidad estática suficiente bajo $F_a=0$.
Se calculó la vida básica por revoluciones con corrección por giro del
anillo exterior\cite{nhbb_vida_exterior_2026}; la mínima es
$1,253\times10^9$ revoluciones del rodillo, como referencia idealizada,
sin atribuir esa vida a los rodamientos adquiridos de marca desconocida.
No se calculan horas sin velocidad manual definida. El montaje directo
del CAD no queda operativo: alojamiento 24 frente a exterior 22\,mm deja
1\,mm radial sin apoyo y exige asiento/adaptador con retención definida.
El apoyo interior debe corresponder al M8 adoptado, con asiento apropiado.

\paragraph{Matriz perforada.}
Se emplea la matriz STEP de 99 mm, 102 perforaciones 6\,mm, espesor perforado
10\,mm y cuerpo central real. El apoyo plano efectivo central es el anillo
de radios 6,5--11,5\,mm. La presión máxima calculada es 3,32\,MPa y el
contacto permanece en compresión en el equilibrio global.
Para evaluar la flexión local se ejecutó FEM tridimensional sobre el STEP,
con cargas equivalentes por canal y fijación adherente de ese anillo,
en tres estados críticos correlacionados.
La malla más fina entrega 36,37\,MPa y 0,01358\,mm.
La envolvente cambia 1,27\% en tensión y 4,05\% en desplazamiento
en el último refinamiento; el anexo detalla los cambios por caso.
La matriz cumple el cribado elástico bajo estas hipótesis, con $n_y=5,77$ respecto del límite de fluencia de referencia.
El material de diseño es AISI 304 con $R_{p0,2}=210$ MPa de referencia para
producto plano H/P a 20 grados Celsius\cite{outokumpu_core_2026}; no es un
certificado de la pieza. La fijación debe transmitir también el torque;
el FEM adherente no demuestra por sí solo que la unión real no patine.

\paragraph{Manivela y funcionamiento.}
La fuerza manual ideal máxima es $2,7804/0,165=16,85$ N.
Con eficiencias supuestas de 0,8 y 0,6 sería 21,06 y 28,08\,N; no son pérdidas
medidas. El tramo recto CAD de 15\,mm presenta flexión nominal 8,39\,MPa.
Faltan el material del brazo y el mecanismo de unión al árbol para verificar
su resistencia completa; un agujero cilíndrico coincidente no garantiza
transmisión de torque. Un atasco no queda limitado al torque de operación:
$T_{\rm mano}=0,165F_{\rm mano}$. No se declara resistencia a un atasco
de magnitud desconocida.

\paragraph{Conclusión de diseño.}
La demanda manual calculada es moderada y los 608 de referencia tienen
capacidad suficiente. Para disponer de un conjunto funcional deben
resolverse tres interfaces concretas: asiento de los 608, apoyo axial del
árbol y transmisión/retención en las uniones matriz--árbol--manivela.
Se requieren además el grado/paso del árbol y el material del brazo.
Estas conclusiones son de cálculo, sin validación de extrusión ni de
funcionamiento experimental. El anexo documenta ecuaciones, FEM y datos
pendientes, sin alterar CAD ni modelos históricos.
